% Academic CV (English).
% Build from the repo root with:  make carolina_branas_cv_english

\documentclass{resume}
% \showframe
\usepackage[left=0.6in,top=0.5in,right=0.6in,bottom=0.45in]{geometry}
\usepackage{vwcol}
%\usepackage{cfr-lm}  % old-style numbers
\usepackage{natbib}
\setlength{\bibsep}{2.5pt}
\usepackage[svgnames]{xcolor}
\usepackage{hyperref}
\usepackage{graphicx}
\hypersetup{%
  colorlinks=true,
  allcolors=DarkRed,
  pdfborderstyle={/S/U/W 1}
}

% compact section spacing
\titlespacing{\section}{0pt}{6pt}{3pt}

\begin{document}

%%%%%%%%%%%%%%%%%%%%%%%%%%%%%%%%%%%%%%%%%%%%%%%%% TITLE
\vspace{-2em}
\begin{minipage}[t]{0.76\textwidth}
    \raggedright
    {\Huge \textbf{Carolina Brañas}}\\[0.3em]
    {\normalsize \textbf{Data Scientist}}\\[0.5em]
    {\footnotesize
    Copenhagen, Denmark $\bullet$
    (+45) 60 90 73 33 $\bullet$
    \href{mailto:carobrasor@gmail.com}{carobrasor@gmail.com} $\bullet$
    \href{https://carobs9.github.io/}{carobs9.github.io} \\[0.2em]
    \href{https://github.com/carobs9}{github.com/carobs9} $\bullet$
    \href{https://www.linkedin.com/in/carolinabranas/}{linkedin.com/in/carolinabranas} $\bullet$
    \href{https://huggingface.co/carobs9}{huggingface.co/carobs9}
    }\\[0.7em]
    {\small
    Data scientist with hands-on experience building machine learning and NLP models,
    designing scalable data workflows, and turning complex data into actionable insights.
    Interests: machine learning, network analysis, natural language processing, geospatial data.
    }
\end{minipage}
\hfill
\begin{minipage}[t]{0.22\textwidth}
    \vspace*{0.2em}
    \begin{flushright}
        \includegraphics[width=2.7cm]{profile.jpg}
    \end{flushright}
\end{minipage}
\vspace{0.2em}

%%%%%%%%%%%%%%%%%%%%%%%%%%%%%%%%%%%%%%%%%%%%%%%%% EXPERIENCE
\section{EXPERIENCE}
\begin{content}

    \begin{position}{Data Science Research Assistant}{Feb 2026 -- Present}{University of Copenhagen, Denmark}{Prof.~Asger Mose Wingender}{}
        \item Research assistant on the Cataloguing Crop Traits and Breeders Across the OECD 1980--2025 project (financed by the \href{https://www.carlsbergfondet.dk/en/}{Carlsberg Foundation}).
        \item Building the first large unified database of all new cereal varieties released in the OECD since 1980 -- their traits, breeding history, and the companies that developed them -- to study how the crop innovation system can raise global food supply without expanding cropland.
        \item Manage a student assistant on the project, coordinating and reviewing their work.
    \end{position}

    \begin{position}{Data Science Research Assistant}{May 2024 -- Dec 2025}{University of Copenhagen, Denmark}{Prof.~Jeanet Bentzen}{}
        \item Contributed to the {\href{https://www.economics.ku.dk/research/externally-funded-research_new/shocking-religion/}{Shocking Religion}} project on religion's economic impact, within a multidisciplinary research team.
        \item Built topic models to uncover thematic trends in text data, and RAG-based LLM pipelines for document insights.
        \item Computed distances from locations to potential earthquake danger areas for each of the four hazard levels (1--4), using Rasterio, GeoPandas and QGIS.
        \item Designed and containerized scalable data workflows; handled large-scale data ingestion and preprocessing; deployed on cloud infrastructure (UCloud).
    \end{position}

    \begin{position}{Data Scientist}{Oct 2023 -- May 2024}{Above Sports, Denmark}{}{}
        \item Developed computer vision models for brand logo detection.
        \item Automated data workflows and worked with product teams to refine output quality.
    \end{position}

    \begin{position}{Marketing Strategist}{Sep 2021 -- May 2022}{Crescendo Collective, United States}{}{}
        \item Analyzed campaign data (Google Analytics, Google Ads) and automated internal reporting with Python.
        \item Ran competitor analysis and audience insights with the data team; presented reports to clients and stakeholders.
    \end{position}

\end{content}

%%%%%%%%%%%%%%%%%%%%%%%%%%%%%%%%%%%%%%%%%%%%%%%%% EDUCATION
\section{EDUCATION}
\begin{content}
    {\bf M.Sc.\ Social Data Science} \hfill {\bf 2022 -- 2024} \\
    University of Copenhagen, Denmark \\
    {\bf \em Thesis:} {\href{https://github.com/carobs9/segregation-madrid}{Mobility and income segregation in Madrid, Spain.}} \\
    {\small {\em Electives:} Advanced Machine Learning for Data Science, Geospatial Data Science, Advanced Network Science (ITU); Natural Language Processing (DIKU).} \\[0.4em]
    {\bf B.Sc.\ Marketing, Computer Science Minor} \hfill {\bf 2018 -- 2022} \\
    Cardinal Stritch University, United States \\
    {\small {\em Honours:} Magna Cum Laude, Dean's List (2018--2022), best graduating GPA of Marketing B.Sc.\ (2022).}
\end{content}

%%%%%%%%%%%%%%%%%%%%%%%%%%%%%%%%%%%%%%%%%%%%%%%%% SKILLS
\section{SKILLS}
\begin{content}
    {\bf Programming \& Data:} Python {\footnotesize (pandas, numpy, pytorch, tensorflow, scikit-learn)}, SQL, Bash, \LaTeX \\
    {\bf Machine Learning \& NLP:} Transformers, topic modeling {\footnotesize (UMAP, HDBSCAN, BERTopic)}, Hugging Face, Weights \& Biases, deep learning, predictive modeling, feature engineering, computer vision \\
    {\bf Analysis, Visualization \& Geospatial:} statistical analysis, data wrangling, Matplotlib, Seaborn, Tableau, Power BI, Geopandas, Rasterio, QGIS \\
    {\bf Cloud \& DevOps:} Docker, Git, Linux, AWS, Azure, UCloud, CI/CD \\
    {\bf Languages:} Spanish {\footnotesize (native)}, Galician {\footnotesize (native)}, English {\footnotesize (professional proficiency)}, Danish {\footnotesize (beginner)}
\end{content}

%%%%%%%%%%%%%%%%%%%%%%%%%%%%%%%%%%%%%%%%%%%%%%%%% AWARDS
\section{AWARDS}
\begin{content}
    \prize[s]{Dean's List}{2018--2022}{{\small every semester of the B.Sc., for a GPA above 3.5}} \\
    \prize[s]{Best Graduating GPA of Marketing B.Sc.}{2022}{{\small highest GPA in Marketing, class of 2022}} \\
    \prize[s]{Academic and Athletic Grant}{2018--2022}{{\small full four-year scholarship for academic excellence and soccer}}
\end{content}

\end{document}
